% Ausblick Kap. 9.2: je Grenze der Arbeit (links) ein Ansatzpunkt für weitere Arbeiten (rechts), mit Bezug auf F/A.
% Eigene Darstellung (Kap. 9, Entwurf 25.11.2026). Inhalte aus Kap. 6.2.3 (Ersatzregler nur Simulation),
%   Kap. 5 (A4 Toleranz, A7 bleibt konzeptionell), Kap. 1.2 (Eingrenzung: keine Real-Navigation, kein Policy-Transfer),
%   Kap. 5.2 (A5 orientiert sich am Rechtsrahmen ohne Erfüllung zu beanspruchen).
% Stil: ba-grafik.tex. Einbinden: % Ausblick Kap. 9.2: je Grenze der Arbeit (links) ein Ansatzpunkt für weitere Arbeiten (rechts), mit Bezug auf F/A.
% Eigene Darstellung (Kap. 9, Entwurf 25.11.2026). Inhalte aus Kap. 6.2.3 (Ersatzregler nur Simulation),
%   Kap. 5 (A4 Toleranz, A7 bleibt konzeptionell), Kap. 1.2 (Eingrenzung: keine Real-Navigation, kein Policy-Transfer),
%   Kap. 5.2 (A5 orientiert sich am Rechtsrahmen ohne Erfüllung zu beanspruchen).
% Stil: ba-grafik.tex. Einbinden: % Ausblick Kap. 9.2: je Grenze der Arbeit (links) ein Ansatzpunkt für weitere Arbeiten (rechts), mit Bezug auf F/A.
% Eigene Darstellung (Kap. 9, Entwurf 25.11.2026). Inhalte aus Kap. 6.2.3 (Ersatzregler nur Simulation),
%   Kap. 5 (A4 Toleranz, A7 bleibt konzeptionell), Kap. 1.2 (Eingrenzung: keine Real-Navigation, kein Policy-Transfer),
%   Kap. 5.2 (A5 orientiert sich am Rechtsrahmen ohne Erfüllung zu beanspruchen).
% Stil: ba-grafik.tex. Einbinden: % Ausblick Kap. 9.2: je Grenze der Arbeit (links) ein Ansatzpunkt für weitere Arbeiten (rechts), mit Bezug auf F/A.
% Eigene Darstellung (Kap. 9, Entwurf 25.11.2026). Inhalte aus Kap. 6.2.3 (Ersatzregler nur Simulation),
%   Kap. 5 (A4 Toleranz, A7 bleibt konzeptionell), Kap. 1.2 (Eingrenzung: keine Real-Navigation, kein Policy-Transfer),
%   Kap. 5.2 (A5 orientiert sich am Rechtsrahmen ohne Erfüllung zu beanspruchen).
% Stil: ba-grafik.tex. Einbinden: \input{Bilder/tikz/ausblick-ansatzpunkte} in einer figure-Umgebung. Nicht skalieren.
\begin{tikzpicture}[ba]
  \tikzset{
    bezug/.style={font=\footnotesize, text=baGrau, anchor=west, text width=11mm, align=left},
    grenze/.style={ba block, anchor=west, align=flush left, text width=62mm, minimum width=66mm, minimum height=13mm, font=\footnotesize},
    ansatz/.style={ba block hervor, anchor=west, align=flush left, text width=62mm, minimum width=66mm, minimum height=13mm, font=\footnotesize},
  }
  % Spalten: Bezug 0,0..1,2 | Grenze 1,4..8,0 | Pfeil 8,0..9,0 | Ansatz 9,0..15,6
  % Reihen 1,3 cm hoch, Raster 1,6 cm (Abstand 0,3 cm), Kopfzeile unten bei 7,8
  \node[font=\small\bfseries, anchor=south west, inner xsep=0pt] at (1.4,7.8) {Grenze dieser Arbeit};
  \node[font=\small\bfseries, anchor=south west, inner xsep=0pt] at (9.0,7.8) {Ansatz für weitere Arbeiten};
  \foreach \y/\f/\g/\a in {
     7.0/{F1, F2}/{Lauf-Policies laufen nur im Ersatzregler der Simulation}/{Transfer einer Lauf-Policy auf den realen G1 über dieselbe Herstellerschnittstelle},
     5.4/{F2, A4}/{Abweichungen zwischen Simulation und Roboter benannt, nicht ausgeglichen}/{Abgleich der Simulation mit Realdaten aus den Messungen dieser Arbeit},
     3.8/{A1}/{Navigation nur in der Simulation, ohne Wahrnehmung der Umgebung}/{Navigation mit Wahrnehmung am realen G1},
     2.2/{F4, A7}/{kein externes Projekt angebunden}/{Folgeprojekt über die Andockpunkte, danach Vergleich zweier Fassungen},
     0.6/{F3, A5}/{Schutzmechanismen ohne Nachweis einer Normerfüllung}/{Sicherheitsnachweis für den Realbetrieb eines balancierenden Humanoiden}%
  }{
    \node[bezug] at (0,\y) {\f};
    \node[grenze] (G) at (1.4,\y) {\g};
    \node[ansatz] (A) at (9.0,\y) {\a};
    \draw[ba pfeil] (G.east) -- (A.west);
  }
\end{tikzpicture}
 in einer figure-Umgebung. Nicht skalieren.
\begin{tikzpicture}[ba]
  \tikzset{
    bezug/.style={font=\footnotesize, text=baGrau, anchor=west, text width=11mm, align=left},
    grenze/.style={ba block, anchor=west, align=flush left, text width=62mm, minimum width=66mm, minimum height=13mm, font=\footnotesize},
    ansatz/.style={ba block hervor, anchor=west, align=flush left, text width=62mm, minimum width=66mm, minimum height=13mm, font=\footnotesize},
  }
  % Spalten: Bezug 0,0..1,2 | Grenze 1,4..8,0 | Pfeil 8,0..9,0 | Ansatz 9,0..15,6
  % Reihen 1,3 cm hoch, Raster 1,6 cm (Abstand 0,3 cm), Kopfzeile unten bei 7,8
  \node[font=\small\bfseries, anchor=south west, inner xsep=0pt] at (1.4,7.8) {Grenze dieser Arbeit};
  \node[font=\small\bfseries, anchor=south west, inner xsep=0pt] at (9.0,7.8) {Ansatz für weitere Arbeiten};
  \foreach \y/\f/\g/\a in {
     7.0/{F1, F2}/{Lauf-Policies laufen nur im Ersatzregler der Simulation}/{Transfer einer Lauf-Policy auf den realen G1 über dieselbe Herstellerschnittstelle},
     5.4/{F2, A4}/{Abweichungen zwischen Simulation und Roboter benannt, nicht ausgeglichen}/{Abgleich der Simulation mit Realdaten aus den Messungen dieser Arbeit},
     3.8/{A1}/{Navigation nur in der Simulation, ohne Wahrnehmung der Umgebung}/{Navigation mit Wahrnehmung am realen G1},
     2.2/{F4, A7}/{kein externes Projekt angebunden}/{Folgeprojekt über die Andockpunkte, danach Vergleich zweier Fassungen},
     0.6/{F3, A5}/{Schutzmechanismen ohne Nachweis einer Normerfüllung}/{Sicherheitsnachweis für den Realbetrieb eines balancierenden Humanoiden}%
  }{
    \node[bezug] at (0,\y) {\f};
    \node[grenze] (G) at (1.4,\y) {\g};
    \node[ansatz] (A) at (9.0,\y) {\a};
    \draw[ba pfeil] (G.east) -- (A.west);
  }
\end{tikzpicture}
 in einer figure-Umgebung. Nicht skalieren.
\begin{tikzpicture}[ba]
  \tikzset{
    bezug/.style={font=\footnotesize, text=baGrau, anchor=west, text width=11mm, align=left},
    grenze/.style={ba block, anchor=west, align=flush left, text width=62mm, minimum width=66mm, minimum height=13mm, font=\footnotesize},
    ansatz/.style={ba block hervor, anchor=west, align=flush left, text width=62mm, minimum width=66mm, minimum height=13mm, font=\footnotesize},
  }
  % Spalten: Bezug 0,0..1,2 | Grenze 1,4..8,0 | Pfeil 8,0..9,0 | Ansatz 9,0..15,6
  % Reihen 1,3 cm hoch, Raster 1,6 cm (Abstand 0,3 cm), Kopfzeile unten bei 7,8
  \node[font=\small\bfseries, anchor=south west, inner xsep=0pt] at (1.4,7.8) {Grenze dieser Arbeit};
  \node[font=\small\bfseries, anchor=south west, inner xsep=0pt] at (9.0,7.8) {Ansatz für weitere Arbeiten};
  \foreach \y/\f/\g/\a in {
     7.0/{F1, F2}/{Lauf-Policies laufen nur im Ersatzregler der Simulation}/{Transfer einer Lauf-Policy auf den realen G1 über dieselbe Herstellerschnittstelle},
     5.4/{F2, A4}/{Abweichungen zwischen Simulation und Roboter benannt, nicht ausgeglichen}/{Abgleich der Simulation mit Realdaten aus den Messungen dieser Arbeit},
     3.8/{A1}/{Navigation nur in der Simulation, ohne Wahrnehmung der Umgebung}/{Navigation mit Wahrnehmung am realen G1},
     2.2/{F4, A7}/{kein externes Projekt angebunden}/{Folgeprojekt über die Andockpunkte, danach Vergleich zweier Fassungen},
     0.6/{F3, A5}/{Schutzmechanismen ohne Nachweis einer Normerfüllung}/{Sicherheitsnachweis für den Realbetrieb eines balancierenden Humanoiden}%
  }{
    \node[bezug] at (0,\y) {\f};
    \node[grenze] (G) at (1.4,\y) {\g};
    \node[ansatz] (A) at (9.0,\y) {\a};
    \draw[ba pfeil] (G.east) -- (A.west);
  }
\end{tikzpicture}
 in einer figure-Umgebung. Nicht skalieren.
\begin{tikzpicture}[ba]
  \tikzset{
    bezug/.style={font=\footnotesize, text=baGrau, anchor=west, text width=11mm, align=left},
    grenze/.style={ba block, anchor=west, align=flush left, text width=62mm, minimum width=66mm, minimum height=13mm, font=\footnotesize},
    ansatz/.style={ba block hervor, anchor=west, align=flush left, text width=62mm, minimum width=66mm, minimum height=13mm, font=\footnotesize},
  }
  % Spalten: Bezug 0,0..1,2 | Grenze 1,4..8,0 | Pfeil 8,0..9,0 | Ansatz 9,0..15,6
  % Reihen 1,3 cm hoch, Raster 1,6 cm (Abstand 0,3 cm), Kopfzeile unten bei 7,8
  \node[font=\small\bfseries, anchor=south west, inner xsep=0pt] at (1.4,7.8) {Grenze dieser Arbeit};
  \node[font=\small\bfseries, anchor=south west, inner xsep=0pt] at (9.0,7.8) {Ansatz für weitere Arbeiten};
  \foreach \y/\f/\g/\a in {
     7.0/{F1, F2}/{Lauf-Policies laufen nur im Ersatzregler der Simulation}/{Transfer einer Lauf-Policy auf den realen G1 über dieselbe Herstellerschnittstelle},
     5.4/{F2, A4}/{Abweichungen zwischen Simulation und Roboter benannt, nicht ausgeglichen}/{Abgleich der Simulation mit Realdaten aus den Messungen dieser Arbeit},
     3.8/{A1}/{Navigation nur in der Simulation, ohne Wahrnehmung der Umgebung}/{Navigation mit Wahrnehmung am realen G1},
     2.2/{F4, A7}/{kein externes Projekt angebunden}/{Folgeprojekt über die Andockpunkte, danach Vergleich zweier Fassungen},
     0.6/{F3, A5}/{Schutzmechanismen ohne Nachweis einer Normerfüllung}/{Sicherheitsnachweis für den Realbetrieb eines balancierenden Humanoiden}%
  }{
    \node[bezug] at (0,\y) {\f};
    \node[grenze] (G) at (1.4,\y) {\g};
    \node[ansatz] (A) at (9.0,\y) {\a};
    \draw[ba pfeil] (G.east) -- (A.west);
  }
\end{tikzpicture}
